\documentclass[a4paper,14pt]{extarticle}

\usepackage[T2A]{fontenc}
\usepackage[utf8]{inputenc}
\usepackage[russian]{babel}
\usepackage{geometry}
\usepackage{setspace}
\usepackage{indentfirst}
\usepackage{array}
\usepackage{booktabs}
\usepackage{listings}
\usepackage{xcolor}

\geometry{
    left=30mm,
    right=15mm,
    top=20mm,
    bottom=20mm
}

\onehalfspacing
\setlength{\parindent}{1.25cm}
\setlength{\parskip}{0pt}
\setlength{\emergencystretch}{3em}

\definecolor{codebg}{RGB}{248,248,248}
\lstdefinestyle{cudastyle}{
    language=C++,
    basicstyle=\ttfamily\normalsize,
    frame=single,
    breaklines=true,
    showstringspaces=false
}

\begin{document}

\begin{titlepage}
\begin{center}
\textbf{МОСКОВСКИЙ АВИАЦИОННЫЙ ИНСТИТУТ}\\
\textbf{(НАЦИОНАЛЬНЫЙ ИССЛЕДОВАТЕЛЬСКИЙ УНИВЕРСИТЕТ)}\\
Институт №8 «Компьютерные науки и прикладная математика»
\end{center}

\vspace{3.5cm}
\begin{center}
\textbf{\Large Лабораторная работа №4}\\
\vspace{0.3cm}
по курсу «Технологии параллельного программирования»\\
\vspace{0.8cm}
\textbf{Работа с матрицами. Метод Гаусса}
\end{center}

\vspace{4cm}
\begin{flushright}
Выполнил: Ахметшин Б.Р.\\
Группа: М8О-403Б-22\\
Преподаватели: А.Ю. Морозов,\\
\hspace*{3.6cm} Е.Е. Заяц
\end{flushright}

\vfill
\begin{center}
Москва, 2026
\end{center}
\end{titlepage}

\section*{Условие}
Цель работы: реализовать метод Гаусса на GPU с выбором главного элемента по столбцу,
используя библиотеку Thrust для поиска максимального элемента, и провести исследование
производительности с помощью профилировщика \texttt{nvprof}.

Выполненная реализация соответствует варианту 3 (решение квадратной СЛАУ
\(Ax=b\)).

\section*{Программное и аппаратное обеспечение}
\textbf{GPU}

\begin{tabular}{>{\raggedright\arraybackslash}p{0.50\textwidth}>{\raggedright\arraybackslash}p{0.34\textwidth}}
Имя & NVIDIA Tesla T4 \\
Compute Capability & 7.5 \\
Объем глобальной памяти & 15637086208 \\
Разделяемая память на блок & 49152 \\
Константная память & 65536 \\
Регистров на блок & 65536 \\
Макс. потоков на блок & (1024, 1024, 64) \\
Макс. размер сетки & (2147483647, 65535, 65535) \\
Количество мультипроцессоров & 40 \\
\end{tabular}

\vspace{0.4cm}
\textbf{CPU}

\begin{tabular}{>{\raggedright\arraybackslash}p{0.50\textwidth}>{\raggedright\arraybackslash}p{0.34\textwidth}}
Архитектура & x86\_64 \\
CPU op-mode(s) & 32-bit, 64-bit \\
Address sizes & 48 bits physical, 48 bits virtual \\
Порядок байт & Little Endian \\
CPU(s) & 12 \\
Имя модели & AMD Ryzen 5 7600 6-Core Processor \\
Ядер / потоков & 6 / 12 \\
\end{tabular}

\vspace{0.4cm}
\textbf{RAM}\\
64 GB

\vspace{0.2cm}
\textbf{ОС}\\
Ubuntu 24.04.4 LTS

\vspace{0.2cm}
\textbf{Компилятор}\\
nvcc, g++

\section*{Метод решения}
Используется классический метод Гаусса с частичным выбором главного элемента
по столбцу:
\begin{enumerate}
    \item на шаге \(k\) выбирается строка с максимальным по модулю элементом
    в текущем столбце (\texttt{thrust::max\_element});
    \item выполняется перестановка строк;
    \item вычисляются коэффициенты исключения;
    \item обнуляются элементы ниже главного в двумерной CUDA-сетке потоков;
    \item после прямого хода выполняется обратный ход (back substitution).
\end{enumerate}

\section*{Описание программы}
\begin{itemize}
    \item \texttt{lab4.cu} --- GPU-реализация: kernels для перестановки строк,
    вычисления коэффициентов и исключения; поиск pivot через Thrust.
    \item \texttt{lab4-cpu.cpp} --- CPU-реализация того же алгоритма для сравнения.
    \item \texttt{generate\_case.py} --- генератор тестовых систем (диагонально доминирующих).
    \item \texttt{run\_case.sh} --- запуск кейса с файла.
    \item \texttt{profile\_nvprof.sh} --- запуск с профилировщиком \texttt{nvprof}.
\end{itemize}

Внутренние замеры:
\begin{itemize}
    \item GPU: \texttt{kernel\_ms}, \texttt{total\_ms} (\texttt{timings\_gpu\_lab4.txt});
    \item CPU: \texttt{total\_ms} (\texttt{timings\_cpu\_lab4.txt});
    \item параметры блоков GPU можно задавать через
    \texttt{LAB4\_SWAP\_THREADS}, \texttt{LAB4\_FACTOR\_THREADS},
    \texttt{LAB4\_ELIM\_BLOCK\_X}, \texttt{LAB4\_ELIM\_BLOCK\_Y}.
\end{itemize}

\textbf{Как из кастомных параметров получается launch-конфигурация kernel'ов:}
\begin{itemize}
    \item \texttt{swapRowsKernel}:\\
    \texttt{block = (LAB4\_SWAP\_THREADS, 1, 1)},\\
    \texttt{grid.x = ceil((n+1)/LAB4\_SWAP\_THREADS)}.
    \item \texttt{computeFactorsKernel}:\\
    \texttt{block = (LAB4\_FACTOR\_THREADS, 1, 1)},\\
    \texttt{grid.x = ceil((n-(k+1))/LAB4\_FACTOR\_THREADS)}.
    \item \texttt{eliminateKernel}:\\
    \texttt{block = (LAB4\_ELIM\_BLOCK\_X, LAB4\_ELIM\_BLOCK\_Y, 1)},\\
    \texttt{grid.x = ceil((n+1-k)/LAB4\_ELIM\_BLOCK\_X)},\\
    \texttt{grid.y = ceil((n-(k+1))/LAB4\_ELIM\_BLOCK\_Y)}.
    \item Поиск главного элемента через \texttt{thrust::max\_element}
    запускает внутренние kernel'ы Thrust/CUB; их конфигурация в \texttt{nvprof}
    отображается отдельно и напрямую через \texttt{LAB4\_*} не задается.
\end{itemize}

\section*{Результаты}
\textbf{Замеры времени:}

\begin{lstlisting}[basicstyle=\ttfamily\scriptsize,frame=single]
+----+--------+---------------+------------+----------------+---------------+
| No | n      | CUDA config   | CPU, ms    | GPU kernel, ms | GPU total, ms |
+----+--------+---------------+------------+----------------+---------------+
| 1  | 128    | swap=128,     | 0.218658   | 794.482422     | 1093.930854   |
|    |        | factor=128,   |            |                |               |
|    |        | elim=(16, 16) |            |                |               |
+----+--------+---------------+------------+----------------+---------------+
| 2  | 128    | swap=256,     |            | 64.523422      | 271.532119    |
|    |        | factor=256,   | 0.218658   |                |               |
|    |        | elim=(32, 32) |            |                |               |
+----+--------+---------------+------------+----------------+---------------+
| 3  | 128    | swap=512,     | 0.218658   | 99.555840      | 296.643681    |
|    |        | factor=512,   |            |                |               |
|    |        | elim=(64, 64) |            |                |               |
+----+--------+---------------+------------+----------------+---------------+
| 4  | 512    | swap=128,     | 13.721943  | 599.301819     | 795.664067    |
|    |        | factor=128,   |            |                |               |
|    |        | elim=(16, 16) |            |                |               |
+----+--------+---------------+------------+----------------+---------------+
| 5  | 512    | swap=256,     | 13.721943  | 514.086426     | 702.878280    |
|    |        | factor=256,   |            |                |               |
|    |        | elim=(32, 32) |            |                |               |
+----+--------+---------------+------------+----------------+---------------+
| 6  | 512    | swap=512,     | 13.721943  | 570.717224     | 780.682874    |
|    |        | factor=512,   |            |                |               |
|    |        | elim=(64, 64) |            |                |               |
+----+--------+---------------+------------+----------------+---------------+
| 7  | 2048   | swap=128,     | 850.001810 | 2534.226562    | 2914.033293   |
|    |        | factor=128,   |            |                |               |
|    |        | elim=(16, 16) |            |                |               |
+----+--------+---------------+------------+----------------+---------------+
| 8  | 2048   | swap=256,     | 850.001810 | 2599.576904    | 2891.329117   |
|    |        | factor=256,   |            |                |               |
|    |        | elim=(32, 32) |            |                |               |
+----+--------+---------------+------------+----------------+---------------+
| 9  | 2048   | swap=512,     | 850.001810 | 2580.171143    | 2901.593170   |
|    |        | factor=512,   |            |                |               |
|    |        | elim=(64, 64) |            |                |               |
+----+--------+---------------+------------+----------------+---------------+
\end{lstlisting}

\newpage
\textbf{Логи профилировщика \texttt{nvprof} (сводка):}

\begin{lstlisting}[basicstyle=\ttfamily\scriptsize,frame=single,breaklines=true]
==10411== NVPROF is profiling process 10411, command: ./lab4
==10411== Profiling application: ./lab4
==10411== Profiling result:
            Type  Time(%)      Time     Calls       Avg       Min       Max  Name
 GPU activities:   36.99%  785.59us       128  6.1370us  4.1600us  7.5520us  [thrust::cuda_cub reduce arg_max kernel]
                   22.05%  468.38us       127  3.6880us  3.4560us  4.3840us  computeFactorsKernel(double const *, double*, int, int, int)
                   20.24%  429.91us       257  1.6720us  1.5680us  12.192us  [CUDA memcpy DtoH]
                   20.08%  426.55us       127  3.3580us  2.8800us  4.4480us  eliminateKernel(double*, double const *, int, int, int)
                    0.64%  13.536us         1  13.536us  13.536us  13.536us  [CUDA memcpy HtoD]

\end{lstlisting}

\textbf{Анализ профилировщика:}
\begin{itemize}
    \item Для профиля \(n=128\) (параметры: swap/factor = 256,
    elim = (16,16)):
    основной вклад дают три группы GPU-активностей:
    \texttt{thrust reduce} (36.99\%), \texttt{computeFactorsKernel} (22.05\%)
    и \texttt{eliminateKernel} (20.08\%).
    Также заметна доля частых \texttt{DtoH} копирований
    малых объемов (20.24\%).
    \item В профиле видно большое число очень коротких запусков
    (порядка \(3\!-\!7\)\,\(\mu s\) на kernel). Для малого размера матрицы это
    приводит к доминированию накладных расходов запуска/синхронизаций.
    \item Копирование крупных массивов не является узким местом:
    HtoD 129\,KB занимает малую долю (0.64\%),
    тогда как
    повторяющиеся мелкие \texttt{DtoH} добавляют значимую задержку.
    \item По таблице замеров видно, что рост \(n\) увеличивает долю полезной
    работы ядра (например, для \(n=2048\) \(kernel\_ms \approx 2.58\!-\!2.60\) c
    при \(total\_ms \approx 2.89\) c), однако реализация все равно уступает CPU
    по общему времени.
\end{itemize}

\section*{Выводы}
По результатам экспериментов и профилирования получены следующие выводы:
\begin{itemize}
    \item Для всех исследованных размеров (\(n=128,512,2048\)) текущая
    GPU-реализация оказалась медленнее CPU по полному времени выполнения.
    Наиболее заметное отставание наблюдается на малых размерах, где
    накладные расходы доминируют.
    \item На \(n=128\) лучшая из протестированных конфигураций
    (swap=256, factor=256, elim=(32,32)) дает
    \(GPU\ total=271.53\) ms против \(CPU=0.22\) ms.
    \item На \(n=512\) лучшая конфигурация
    (swap=256, factor=256,
    elim=(32,32)) дает
    \(GPU\ total=702.88\) ms против \(CPU=13.72\) ms.
    \item На \(n=2048\) лучшая конфигурация
    (swap=256, factor=256,
    elim=(32,32)) дает
    \(GPU\ total=2891.33\) ms против \(CPU=850.00\) ms.
    \item Профиль \texttt{nvprof} подтверждает, что существенную долю времени
    занимают многократные короткие запуски kernel'ов и частые мелкие
    синхронизирующие копирования \texttt{DtoH}. Это объясняет, почему даже при
    корректной двумерной сетке потоков GPU не показывает ускорения в текущей
    постановке.
    \item Для повышения эффективности целесообразно уменьшать число запусков
    kernel'ов и CPU-GPU синхронизаций (укрупнение шагов, сокращение мелких
    \texttt{DtoH}, перенос большего объема логики выбора pivot на GPU).
\end{itemize}

\end{document}
